% !TeX root = ../nc-del-pezzo.tex
% 2026-11-20 JST / v0.13
\section{全同値群oidと可換零切断近傍の有限型モジュライ}
\label{sec:finite-moduli}
\label{sec:source-near-commutative-zero}

定理\ref{thm:simultaneous-smooth-ncdp-open}，
系\ref{cor:finite-equivalence-unramified}，
命題\ref{prop:simultaneous-hochschild-rigid-open}の
共通の $G_d$ 不変開を，全係数を含む有限型開集合 $V$ とし，
定理\ref{thm:finite-anticanonical-isom}の全同値群oidを
\[
 \mathcal E\underset{t}{\overset{s}{\rightrightarrows}}V
\]
と書く。ここでは $V$ を一つの横断sliceに制限していない。
必要なら滑らかな古典係数の開集合 $U_d$ と方向錐 $C_{n,d}$ の積に制限する。
この制限は滑らかな可換零切断
$Z=U_d\times\{0\}\subset V$ の全体を保つ。
以下の相対次元は，この全係数表示の座標変換自由度を数えるものである。

\subsection{有効な座標変換と形式横断slice}

第\ref{sec:definitions}節の群 $G_d$ について
\[
 G_1\simeq\operatorname{GL}_n,
 \qquad
 G_2\simeq\operatorname{GL}_{n+1}/\mu_2,
 \qquad
 g_d:=\dim G_d=
 \begin{cases}n^2,&d=1,\\(n+1)^2,&d=2.\end{cases}
\]
最初の同型は $(A,t)\mapsto t^{-1}A$ による。
次数2では $A\mapsto[(A,1)]$ が核 $\mu_2$ のfppf全射であり，
基底上で第二成分の平方根を取れば全射性を確認できる。
従って両群は滑らかである。

有限ambientと選んだ中心元は
定理\ref{thm:simultaneous-degree-one-central-family}，
定理\ref{thm:simultaneous-degree-two-central-family}により
$\operatorname{GL}$ 同変である。
次数2の $-I$ は $x_i\mapsto-x_i$，$z\mapsto z$ と作用し，
これは次数スカラー $\varepsilon_{-1}$ である。
そのtails上の作用は，graded加群の次数 $j$ に $(-1)^j$ を掛ける自然同型によって，
構造対象の標識を保つ恒等関手に同一視される。
この自然同型は積と両立するから，有効群 $G_d$ による作用と射
\begin{equation}\label{eq:coordinate-action-to-equivalences}
 \rho:G_d\times V\longrightarrow\mathcal E,
 \qquad(g,a)\longmapsto[a\xrightarrow{\ g\ }ga]
\end{equation}
が定まる。$s\rho=\operatorname{pr}_V$，$t\rho(g,a)=ga$ である。

\begin{lemma}[全パラメータ空間の形式的な座標自由度]
\label{lem:full-coefficient-formal-slice}
$x=(f,0)\in Z(\mathbb C)$ とし，その多様体を $X$ とする。
全framed Abel変形環を $R_X$ と書く。
全係数の同じ族が定めるframed変形への自然射について，
非標準的だが両立する同型
\begin{equation}\label{eq:full-coefficient-formal-product}
 \widehat{\mathcal O}_{V,x}
       \simeq R_X[[u_1,\ldots,u_{g_d}]]
\end{equation}
がある。その自然射は右辺の $R_X$ の包含に対応する。
\end{lemma}
\begin{proof}
係数空間の接空間で
$\operatorname{Lie}(G_d)\cdot f$ の補空間 $N$ を取り，
$S=(f+N)\cap U_d$ と置く。
$H^0(X,T_X)=0$ と正規形の商表示
\cite[定理8.8]{PoissonDP}により，軌道接写像は単射であり，
$\dim N=b_{n,d}$ である。
従って $G_d\times S\to U_d$ は $(1,f)$ で\'{e}taleである。
有限自己同型群の商を取っているのではなく，群の恒等元の形式近傍を使っている。

錐を含む写像
\[
 \Phi:G_d\times S\times C_{n,d}\longrightarrow U_d\times C_{n,d},
 \qquad(g,h,c)\longmapsto(gh,gc)
\]
は，$(g,h,c)\mapsto(g,h,gc)$ という自己同型と，
$G_d\times S\to U_d$ と錐の恒等との積の合成である。
従って $\Phi$ は $(1,f,0)$ で\'{e}taleであり，
次数2の錐が特異でも
\[
 \widehat V_x\simeq
     \widehat{G_d}_1\times\widehat{(S\times C_{n,d})}_{(f,0)}
\]
となる。開条件はこの完成を変えない。

任意の局所Artin環 $A$ に対して，$x$ に還元するパラメータは
$a=g\cdot q$，$g\equiv1$，$q\in S\times C_{n,d}$ と一意に書ける。
座標変換による同値は特殊ファイバー上で恒等に還元するため，
$a$ と $q$ は同じframed Abel変形を与える。
定理\ref{thm:simultaneous-central-all-artin}は，この同じ有限中心商族の
slice上の完備化が全framed変形関手を表すことを与える。
従って変形関手への射は $q$ への射影であり，
$\widehat{G_d}_1\simeq\operatorname{Spf}\mathbb C[[u_1,\ldots,u_{g_d}]]$
と合わせて結論を得る。
\end{proof}

\begin{remark}[classifying mapに沿う，strictな同変性を使わない証明]
\label{rem:formal-slice-without-strict-equivariance}
補題\ref{lem:full-coefficient-formal-slice}の環の同定は，
有限中心元のstrictな $G_d$-同変性に依存せず確認できる。
係数空間の線型補空間を選べば，まず環だけの同型
\[
 A:=\widehat{\mathcal O}_{V,x}\simeq C_X[[u_1,\ldots,u_{g_d}]],
 \qquad C_X:=\widehat{\mathcal O}_{S\times C_{n,d},(f,0)}
\]
を得る。sliceへの制限は $u=0$ であり，
族のclassifying mapを $q^*:R_X\to A$ と書く。
定理\ref{thm:simultaneous-central-all-artin}によれば
$R_X\to A\to C_X$ は同型である。
その同型で $R_X$ と $C_X$ を同一視した後，
$C_X[[u]]\to A$ を，$C_X$ 上では $q^*$，各 $u_i$ 上では選んだ係数座標で定める。
これは $u=0$ を法として恒等であり，余接空間上では対角ブロックが恒等の三角行列となる。
従って補題\ref{lem:formal-endomorphism-inverse}により同型である。
この証明はPl\"ucker錐の特異点でも有効であり，
単なる抽象的な環同型でなく，実際のclassifying mapに沿う形式的滑らかさを与える。
\end{remark}

\subsection{可換点を結ぶ任意の矢の完成}

\begin{theorem}[可換零切断でのsource・targetの完成]
\label{thm:equivalence-source-formal-smooth}
$x,y\in Z(\mathbb C)$ とし，$\alpha:x\to y$ を全構造対象付き同値の任意の矢とする。
両端のframed変形環を $\alpha$ によって同一視して $R_X$ と書けば，
\begin{align}
 \widehat{\mathcal O}_{V,x}
    &\simeq R_X[[u_1,\ldots,u_{g_d}]],\notag\\
 \widehat{\mathcal O}_{V,y}
    &\simeq R_X[[v_1,\ldots,v_{g_d}]],\notag\\
 \widehat{\mathcal O}_{\mathcal E,\alpha}
    &\simeq R_X[[u_1,\ldots,u_{g_d},v_1,\ldots,v_{g_d}]].
 \label{eq:equivalence-completed-local-ring}
\end{align}
$s^*,t^*$ はそれぞれ対応する部分形式変数環の包含である。
従って $s,t$ は $\alpha$ で形式的に滑らかであり，
実際に $\alpha$ の開近傍で滑らか，相対次元は $g_d$ である。

さらに，$\alpha$ を表す正規形座標変換 $g_0\in G_d(\mathbb C)$ を取ると，
\eqref{eq:coordinate-action-to-equivalences}は
$(g_0,x)$ と $\alpha$ における完備局所環の同型を誘導し，
$(g_0,x)$ で\'{e}taleである。
\end{theorem}
\begin{proof}
可換点での全 $\mathcal O$ 保存Abel同値は多様体の同型である
（定理\ref{thm:unframed-formal-quotient}の証明）。
正規形の商表示によって，$\alpha$ はある $g_0$ から誘導される。
枠の変更により，$y$ の変形関手を $X$ のframed変形関手へ移す。

局所Artin環 $A$ 上の $\alpha$ の持上げは，
$x,y$ のパラメータ持上げ $a,b$ と，還元が $\alpha$ となる同値の組である。
枠を変更した後では，二つのframed変形類が等しいときに限りその同値が存在する。
また定理\ref{thm:quantization-framed-groupoid}と
定理\ref{thm:simultaneous-central-all-artin}により，存在するときは一意である。
従って矢の形式関手は文字通り
\[
 \widehat V_x\times_{\operatorname{Spf}R_X}\widehat V_y
\]
である。有限自己同型群の他の元は別の閉矢であり，
固定した $\alpha$ における完成へ混入しない。
補題\ref{lem:full-coefficient-formal-slice}を両端に適用すれば
\eqref{eq:equivalence-completed-local-ring}を得る。
従ってsource・targetの局所環写像は，Artin小拡大に対する持上げ判定を満たす。
有限型射の局所判定\cite[Lemma~15.39.6, Tag~07VH]{StacksProject}により，
$s,t$ は $\alpha$ で滑らかであり，相対次元は加えた形式変数の数 $g_d$ である。
ここまでの証明は注意\ref{rem:formal-slice-without-strict-equivariance}を用いれば，
有限族のstrictな群同変性を仮定する必要がない。

座標作用との一致も明示できる。
$a=g_aq_a$，$g_0^{-1}b=g_bq_b$ を上の一意な形式正規化とし，
$g_a,g_b\equiv1$ とする。
与えられた矢をこの座標変換でsliceへ移すと，
特殊ファイバー上で恒等となる $q_a$ と $q_b$ のframed同値である。
従って $q_a=q_b$ であり，その矢は一意だから，元の矢は
\[
       g=g_0g_bg_a^{-1}
\]
の座標作用と等しい。
逆に，$g\equiv g_0$ と $a\equiv x$ はこのような矢を定める。
一意性は形式正規化の一意性から従う。
これにより $\rho$ は完成上の同型を与える。

有限型複素スキームの射について，この完成局所環の同型は
\'{e}tale性を与える\cite[Lemma~41.11.3, Tag~039M]{StacksProject}。
これは相対次元の別の確認も与える。
実際，$s\rho$ は滑らかな群 $G_d$ からの射影であり，
$t\rho$ も $(g,a)\mapsto(g^{-1},ga)$ によって同じ射影へ移る。
\end{proof}

\begin{remark}[相対滑らかさと絶対特異点]
この完成は具体的に
\[
 \widehat{\mathcal O}_{\mathcal E,\alpha}\simeq
 \begin{cases}
 \mathbb C[[t_1,\ldots,t_{b_{n,1}},D_1,\ldots,D_n,
                u_1,\ldots,u_{n^2},v_1,\ldots,v_{n^2}]],&d=1,\\[2pt]
 \mathbb C[[t_1,\ldots,t_{b_{n,2}},p_{ij},
                u_1,\ldots,u_{(n+1)^2},v_1,\ldots,v_{(n+1)^2}]]/I_{\mathrm{Pl}},&d=2.
 \end{cases}
\]
次数2・$n\ge3$ のPl\"ucker特異点は残っている。
sourceが滑らかであるとは，この特異な基底に対して自由な座標変換変数を加えることを意味する。
また $\mathcal E\to V\times V$ がunramifiedであることと，
$s,t$ の正の相対次元 $g_d$ は矛盾しない。
\end{remark}

\subsection{片端が可換である場合}

\begin{lemma}[可換零切断の幾何学的飽和性]
\label{lem:commutative-zero-equivalence-saturated}
任意の幾何学的矢 $\alpha:x\to y$ について，$x\in Z$ ならば $y\in Z$ である。
\end{lemma}
\begin{proof}
反標準認識\eqref{eq:finite-anticanonical-isom}により，
$B_y^{(r)}$ は可換環 $B_x^{(r)}$ と次数付き同型であり，従って可換である。
$n=2$ なら $r=1$ なので $B_y$ 自体が可換である。
$n\ge3$ とし，$y$ の方向が非零だと仮定する。
命題\ref{lem:simultaneous-central-flatness}により，
中央一次元 $u\in\mathcal K_y$ を一つ選べば $u$ は $B_y$ 上の非零因子である。
任意の斉次元 $a,b$ に対し，$u^ia,u^jb$ の次数が $r$ の倍数となるように
$i,j\ge0$ を選ぶ。すると
\[
       0=[u^ia,u^jb]=u^{i+j}[a,b]
\]
であるから $[a,b]=0$，従って $B_y$ は可換となる。

一方，非零方向では\eqref{eq:simultaneous-degree-one-relations}または
\eqref{eq:simultaneous-ambient-relations}のいずれかに
\[
 [x_i,y]=D_i z\ne0\quad(d=1),\qquad
 [x_i,x_j]=p_{ij}z\ne0\quad(d=2)
\]
が現れる。ここで $z\ne0$ はambientのPBW表示から従い，
中央関係式の次数 $6,4$ はそれぞれ $\deg z=3,2$ より大きいので，商 $B_y$ でも保たれる。
矛盾により方向は零である。
従って $y$ は可換零切断にある。
この議論は低次数中心の最小性を仮定せず，中央一次非零因子の存在だけを使っている。
\end{proof}

\begin{corollary}\label{cor:source-smooth-over-zero-section}
$s$ または $t$ の像が $Z$ に属する全ての点で，$s,t$ はともに滑らかであり，
相対次元は $g_d$ である。
\end{corollary}
\begin{proof}
前補題により，片端が零切断なら両端が零切断である。
複素閉矢では定理\ref{thm:equivalence-source-formal-smooth}を適用できる。
滑らかでないlocusは閉集合であり，有限型複素スキームの非空な閉部分は複素閉点を持つ。
従って零切断上の全ての点で同じ結論となる。
\end{proof}

\begin{corollary}[零切断のschemeとしての保存]
\label{cor:zero-section-scheme-saturated}
$Z$ のイデアル層を $\mathcal J$ とすると，$\mathcal E$ 上で
\[
       s^{-1}\mathcal J\cdot\mathcal O_{\mathcal E}
        =t^{-1}\mathcal J\cdot\mathcal O_{\mathcal E}
\]
が成り立つ。従って任意の基底上の矢について，sourceが $Z$ を通って因子化することと，
targetが $Z$ を通って因子化することは同値である。
\end{corollary}
\begin{proof}
零切断に接しない点では，前補題により両イデアルとも単位イデアルである。
零切断に接する複素閉矢 $\alpha$ を固定し，
$A=\widehat{\mathcal O}_{\mathcal E,\alpha}$ とする。
その任意の局所Artin商 $T$ 上の矢を $a\to b$ と書く。
古典同型は基本偏極 $H$ を保ち，同じ有限族のArtin基底変更は
平坦Abel変形，豊富性，section飽和性を満たす
（定理\ref{thm:simultaneous-relative-tails}と
定理\ref{thm:simultaneous-central-all-artin}）。
従って命題\ref{prop:artin-full-graded-recognition}が適用でき，
この矢は次数スカラーを法とした通常のgraded環同型 $B_a\simeq B_b$ により表される。

$a$ が零sectionを通るなら $B_a$ は可換であり，従って $B_b$ も可換である。
$(B_b)_{\deg z}$ は有限自由 $T$-加群で，$z$ の剰余体への還元は零でない。
従って $Tz$ はこの加群の自由な直和因子である。
表示関係式 $D_i z=0$ または $p_{ij}z=0$ から全方向パラメータが零となり，
$b$ も零sectionを通る。逆向きも同様である。

$A$ の極大イデアルを $\mathfrak m$，両方向イデアルを $J_s,J_t$ とする。
上の結論を全ての $A/(J_s+\mathfrak m^k)$ に適用すれば
$J_t\subset\bigcap_k(J_s+\mathfrak m^k)=J_s$ を得る。
最後の等号は完備Noether局所環のKrull交叉による。
逆も同じなので，完備局所環で両イデアルは等しい。
完備化はNoether局所環上忠実平坦であり，両イデアルはcoherentなので，
全複素閉点でのこの一致はイデアル層自身の一致を与える。
この証明は有限中心元のstrictな群同変性を用いない。
\end{proof}

\begin{proposition}[単位での滑らかなsourceを持つ基底開集合]
\label{prop:unit-source-smooth-open}
$\operatorname{Sm}_{g_d}(s)\subset\mathcal E$ を，$s$ が滑らかで相対次元 $g_d$ である点の開集合とする。
単位切断を $e:V\to\mathcal E$ とすると
\[
 V^{\mathrm{unit}}:=e^{-1}\bigl(\operatorname{Sm}_{g_d}(s)\bigr)
\]
は $Z$ 全体を含む開集合である。
逆射は単位を固定してsourceとtargetを交換するため，
$s$ の代わりに $t$ を用いても同じ開集合となる。
従って $V^{\mathrm{unit}}$ への制限後は，全ての単位でsourceとtargetが
相対次元 $g_d$ の滑らかな射である。
この開条件は座標変換と両立する。
\end{proposition}
\begin{proof}
滑らかな射の相対次元は局所定数なので，$\operatorname{Sm}_{g_d}(s)$ は開である。
零切断上の単位には定理\ref{thm:equivalence-source-formal-smooth}を適用する。
最後の主張は，座標変換がこの群oidのsource・targetと単位を保つ自己同型を与えることによる。
\end{proof}

\begin{remark}[単位近傍の滑らかさと全矢の制御]
以上は，$\mathcal E$ の中のsourceとtargetの共通smooth locusが，
零切断に接する全ての矢を含むことを証明する。
この議論だけでは，$Z$ を含む一つの基底開集合 $W\subset V$ を選び，
$\mathcal E_W$ の全ての矢でsource・targetを滑らかにできたとは結論しない。
非smoothな矢の閉集合が零切断に接しなくても，そのsource像の閉包は零切断へ達し得る。
ここでsourceのproper性は仮定していないため，その像を閉集合として除くことはできない。
全矢の制御には後の定理\ref{thm:full-source-smooth-zero-neighborhood}の
Casimir Jacobianを用い，その後に商の代数性を証明する。
全モジュライ内での開性は定理\ref{thm:zero-neighborhood-open-full-moduli}による別の結論であり，
全非可換対象の被覆は未解決である。

単位での滑らかさだけから全矢の滑らかさを結論することもできない。
例えば節点曲線の正規化 $\widetilde C\to C$ のkernel群oid
$\widetilde C\times_C\widetilde C\rightrightarrows\widetilde C$ を考える。
節点の同じ枝を二度取る単位の完成は対角なので，sourceはそこで同型である。
異なる二つの枝を結ぶ矢の完備局所環は $\mathbb C$ となり，
そのsourceは曲線上の閉点の包含であって滑らかではない。
必要なら滑らかな群の自由度を積に加え，同じ障害を正の単位相対次元でも作れる。
従って命題\ref{prop:unit-source-smooth-open}は全矢の制御を始める開条件であり，
それ自身を滑らかなpresentationの完成とみなしてはいけない。
\end{remark}

\subsection{基本シフトを保つ矢の開部分群oid}

以下では有限型開集合 $V$ を
命題\ref{prop:simultaneous-hochschild-rigid-open}の開集合まで縮める。
これは滑らかな可換模型の零切断全体を保つ。
$V\times V$ 上で，通常の次数付き同型と全構造対象付き同値の層を
\[
 \mathcal H=
 \underline{\operatorname{Isom}}_{\mathrm{gr}}(B_1,B_2)/\mathbb G_m,
 \qquad
 \mathcal E=
 \underline{\operatorname{Isom}}_{\mathrm{gr}}
       (B_1^{(r)},B_2^{(r)})/\mathbb G_m,
 \qquad r=n-1
\]
と書く。両者は補題\ref{lem:graded-isom-affine-quotient}により
有限表示アフィンスキームであり，後者は
定理\ref{thm:finite-anticanonical-isom}の全同値層である。
Veroneseへの制限は自然な射 $j:\mathcal H\to\mathcal E$ を与える。

\begin{lemma}[通常次数付き表示の忘却の核]
\label{lem:ordinary-graded-recognition-injective}
$j$ はmonomorphismである。
この主張には $HH^1$ 消滅や低次数中心条件は不要である。
\end{lemma}
\begin{proof}
任意の許された基底上でfppf局所的に二つの次数付き同型を代表で取り，
その合成を $B$ の次数付き自己同型 $\rho$ に帰着する。
$j(\rho)$ が恒等なら，誘導されるtails自己同値は標識付き恒等関手へ自然同型である。
全シフトを $s$ とし，$E_i=s^{-i}\mathcal O$ と置く。
飽和性から
\[
 \operatorname{End}(E_i)=\mathcal O_T,
 \qquad \operatorname{Hom}(E_{i+1},E_i)=B_1
\]
であり，$B_1$ は正階数局所自由で忠実である。
次数付き同型はシフトと厳密に可換するので，
この自然同型の $E_i$ 上の成分を単元 $\lambda_i$ と書ける。
構造対象の標識により $\lambda_0=1$ である。
隣接Homの自然性と，各添字で同じ $\rho|_{B_1}$ が作用することから，
$\lambda_{i+1}/\lambda_i=c$ は $i$ によらない。
一般の $\operatorname{Hom}(E_{i+m},E_i)=B_m$ に自然性を適用すると
$\rho|_{B_m}=c^m\operatorname{Id}$ となる。
従って核は既に割った次数スカラーだけであり，$j$ は層のmonomorphismである。
両辺は有限表示スキームなので，これは任意の基底上のmonomorphismも意味する。
\end{proof}

\begin{theorem}[Hochschild剛性による基本シフト認識の開性]
\label{thm:ordinary-graded-arrows-open}
上で選んだ $V$ に対して，$j:\mathcal H\to\mathcal E$ は開埋込みである。
その像は全ての単位を含み，合成と逆射で保たれる開部分群oidである。
従って，全 $\mathcal O$ 保存同値の中で基本シフトと可換するという条件は，
矢のスキーム上のZariski開条件となる。
この像は可換零切断の二点を結ぶ全ての矢を含む。
$n\ge3$ の低次数中心の開集合 $U$ では
$\mathcal H|_{U\times U}=\mathcal E|_{U\times U}$ であり，
$n=2$ では $V$ 全体で一致する。
\end{theorem}
\begin{proof}
前補題により $j$ はmonomorphismである。
その始域と終域は有限型 $\mathbb C$-スキームだから，$j$ は有限表示である。
閉点におけるArtin持上げ条件を確認する。

$R'\twoheadrightarrow R$ を剰余体 $\mathbb C$ の局所Artin環の小拡大とし，
$\mathcal E(R')$ の矢 $F$ と，その還元の $\mathcal H(R)$ への持上げを固定する。
局所Artin環の可逆加群は自明なので，必要な $\mathbb G_m$-torsorを自明化し，
これらを実際の次数付き環同型と標識付きtails同値で代表させられる。
さらに $R^\times$ の元は $(R')^\times$ へ持ち上がるので，
代表の還元も所望のものに合わせられる。
対象側の基本シフトを $s_1,s_2$ とし，
\[
 G=s_2^{-1}\circ F\circ s_1\circ F^{-1}
\]
という第二の圏の自己同値を考える。
与えられた $\mathcal H(R)$ の矢は，$G_R$ と恒等関手との指定された自然同型を与える。

\emph{可逆bimoduleがperfectである理由。}
完全Abel同値は複体へ項ごとに作用する標準dg導来同値を誘導し，
直和とcompact対象を保つ。
有限窓のdg模型を $C_{R'}$ とすれば，$G$ は可逆な
$C_{R'}$-bimodule $P$ によって表される。
その準逆bimoduleを $Q$ とすると，
\[
 P\otimes^L_{C_{R'}}Q\simeq C_{R'},
 \qquad Q\otimes^L_{C_{R'}}P\simeq C_{R'}
\]
である。
bimoduleの圏で一方の作用に関して $P$ とtensorする関手は，
$Q$ とのtensorを準逆とする，直和を保つ導来同値である。
従ってcompact対象をcompact対象へ送る。
$C_{R'}$ はsmoothなので対角は $C_{R'}^{\mathrm e}$ 上perfectであり，
その像である $P$ も $C_{R'}^{\mathrm e}$ 上perfectとなる。
ここでは一般のAbel変形のdg模型を仮定せず，有限族から来る既存の模型だけを使った。

\emph{指定された比較同型の持上げ。}
$P$ の特殊ファイバーは標準対角である。
系\ref{cor:noncommutative-diagonal-artin-rigidity}から
$P\simeq C_{R'}$ を得る。
その $R$ 上への還元が与えられた $G_R\simeq\operatorname{Id}$ と異なる場合も，
差は対角の自己同型である。
命題\ref{prop:simultaneous-hochschild-rigid-open}により
\[
 \operatorname{End}_{D(C_R^{\mathrm e})}(C_R)=HH_R^0(C_R)=R
\]
なので，その差は一つの単元 $c\in R^\times$ に等しい。
$c$ を $(R')^\times$ へ持ち上げて同型をスカラー倍すれば，
指定された $R$ 上の比較と厳密に両立する同型を取れる。
bimoduleの同型は導来関手の自然同型を与え，
そのheartへの制限により $Fs_1\simeq s_2F$ を得る。

この比較を正負に整合的に反復し，構造対象の標識と組み合わせる。
全シフトのsection環の飽和性と豊富対による復元から，
$F$ は通常の次数付き同型 $B_1\simeq B_2$ によって誘導される。
指定された比較を持ち上げたため，その還元は所望の $\mathcal H(R)$ の矢である。
同じ $\mathcal E$ の矢を持ち上げる二つの類は前補題により一致する。
したがって $j$ は全閉点で形式的\'{e}taleである。

局所Noetherな有限型射のArtin持上げ判定
\cite[Tag 02HY]{StacksProject}から $j$ は全複素閉点でsmoothである。
非smooth locusは閉であり，有限型複素スキームの全閉点を避けるので空である。
従って $j$ はsmoothであり，
monomorphismなので相対次元0，すなわち\'{e}taleである。
\'{E}tale monomorphismは開埋込みである
\cite[Tag 025G]{StacksProject}。
合成・逆射・単位との両立は通常の次数付き同型から明らかである。
可換点の矢は基本偏極を保つので，通常の切断環へ持ち上がる。
最後の二主張は系\ref{cor:finite-full-graded-recognition}と $r=1$ の場合による。
\end{proof}

\begin{remark}[矢の開性と対象の開近傍は異なる]
この定理は $\mathcal E$ 内の開部分群oidを与える。
それだけから，零切断全体を含む対象側の開集合 $W\subset V$ を取り
$\mathcal H|_{W\times W}=\mathcal E|_{W\times W}$ とできるとは結論しない。
また $j$ の\'{e}tale性は，群oidのsource・targetの\'{e}tale性やsmoothnessとは別の主張である。
\end{remark}

\begin{corollary}[可換零切断近傍における通常次数付き表示の全同値認識]
\label{cor:zero-neighborhood-full-graded-recognition}
$V_d$ を命題\ref{prop:simultaneous-hochschild-rigid-open}の開集合内に取る。
$n\ge3$ では定理\ref{thm:low-center-zero-neighborhood}の $q=r=n-1$ に対する開集合を $W_d=W_{d,r}$ とし，
$n=2$ では $W_d=V_d$ とする。
いずれも滑らかな可換零切断 $Z_d$ 全体を含み，
任意の許された基底上の二点 $a,b:T\to W_d$ に対して
\begin{equation}\label{eq:zero-neighborhood-full-graded-recognition}
 \underline{\operatorname{Eq}}_{\mathcal O}(\QGr B_a,\QGr B_b)
 \simeq
 \underline{\operatorname{Isom}}_{\mathrm{gr},T}(B_a,B_b)/\mathbb G_m
\end{equation}
が自然に成立する。
特に二点の方向が基底上で冪零となる場合や，零方向と非零方向の
特殊化を含む場合にも，この同値層の認識は成立する。
さらに通常の次数付き同型から反標準Veroneseの同型への制限は，
$W_d\times W_d$ 上fppf $\mu_r$-torsorである。
\end{corollary}
\begin{proof}
$n=2$ では $r=1$ なので定理\ref{thm:finite-anticanonical-isom}そのものである。
$n\ge3$ とする。
通常の次数付き同型と全同値のアフィン有限表示スキームを
\[
 \mathcal H=
  \underline{\operatorname{Isom}}_{\mathrm{gr}}(B_1,B_2)/\mathbb G_m,
 \qquad
 \mathcal E=
  \underline{\operatorname{Isom}}_{\mathrm{gr}}(B_1^{(r)},B_2^{(r)})/\mathbb G_m
\]
と書く。
定理\ref{thm:ordinary-graded-arrows-open}により，自然な射
$j:\mathcal H\to\mathcal E$ は開埋込みである。
その $W_d\times W_d$ 上の制限が全射であることを示す。

有限型 $\mathbb C$-スキーム上の開埋込みなので，閉点について確認すれば十分である。
$\mathcal E$ の一つの閉点を，$a,b\in W_d(\mathbb C)$ を結ぶ同値と考える。
補題\ref{lem:commutative-zero-equivalence-saturated}により，
$a,b$ の一方が零切断に属すれば他方も属す。
この場合，定理\ref{thm:ordinary-graded-arrows-open}の可換点の主張により
この矢は $\mathcal H$ に入る。
両方が非零方向なら，定理\ref{thm:low-center-zero-neighborhood}により $a,b\in U_{d,r}$ なので，
系\ref{cor:finite-full-graded-recognition}からやはり $\mathcal H$ に入る。
従って開埋込み $j$ の像の補集合には閉点がなく，補集合は空である。
全射な開埋込みはスキームの同型だから，
任意の非被約基底変更後も $\mathcal H\simeq\mathcal E$ であり，
式\eqref{eq:zero-neighborhood-full-graded-recognition}を得る。
ここでは零切断上の低次数中心の最小性を仮定していない。
また零切断のschemeとしてのsaturationは不要であり，
幾何学的saturationのみを用いた。

最後に両側の未商化Isomスキームは，この共通の商層上の
$\mathbb G_m$-torsorである。
制限写像はスカラーに関して $\lambda\mapsto\lambda^r$ と両立する。
fppf局所的に両torsorを自明化すれば，この射は
$\mathbb G_m\to\mathbb G_m$，$\lambda\mapsto c\lambda^r$
の形になる。これは $\mu_r$-torsorなので主張が従う。
\end{proof}

\begin{remark}[結論の範囲]
以上で，滑らかな可換零切断の一つの開近傍において，
基本シフトを標識に追加せず全 $\mathcal O$ 保存同値を通常の $B$ から認識できる。
通常の次数付き同型の定義は，重みの大きい生成元へ低い重みの生成元の
同次数多項式を加える変換も許す。
実際の全同値についてこれらの付加項が消えることは，
後の定理\ref{thm:zero-neighborhood-finite-diagonal}の準備で証明する。
全同値群oidのsource・targetのsmoothnessには，次項以降の中心Jacobianによる
独立した持上げの証明を使う。
その商の全モジュライ内での開性には
定理\ref{thm:zero-neighborhood-open-full-moduli}の独立した比較を用いる。
全nc対象を被覆するかどうかは未解決である。
方向解消のproper性は悪い閉集合の像を閉にするために用いており，
その射を平坦降下の被覆として使ってはいない。
\end{remark}

\subsection{方向解消上の規格化中心とCasimir係数の微分}

本項では，選択した有限中心元のCasimir係数方向の変化を調べる。
中心元の選択が係数の線型平行移動と厳密に可換することは仮定しない。
必要なのは，その変化の微分が可換零切断で恒等となることである。

中心元の内部次数を $e_1=6$，$e_2=4$ とする。
滑らかな可換係数の開集合を $U_d$ とし，
既存の同時中心族の定義域を $U_d\times C_{n,d}$ に制限する。
さらに，各 $m\le e_d$ の完全対称化
\[
 (\operatorname{Sym}\mathcal V)_m\longrightarrow A_m
\]
が同型となる有限個の開条件を課す。
両辺は同じ階数の有限局所自由加群であり，零切断では同型なので，
これは可換零切断全体を含む不変開集合である。
ここで $\mathcal V$ は重み付き生成空間であり，完全対称化には $y,z$ も含める。
この段階の開集合を $V$，その零切断を $Z$ と書く。
選択した中心元を $\Omega$ とし，完全対称化座標の $z^2$ 係数を読む写像を
\[
       \lambda:A_{e_d}\longrightarrow\mathcal O_V
\]
と書けば，既存の規格化は $\lambda(\Omega)=1/2$ である。

方向解消を
\[
 q:Y\longrightarrow V
\]
と書く。次数1では $\mathbb P(E^*)$ 上の $\mathcal O(-1)$，
次数2では $\operatorname{Gr}(2,E^*)$ 上の $\mathcal O(-1)$ の全空間に
係数空間を掛け，$V$ へ引き戻したものである。
$q$ はproperで，零方向の外では同型である。
例外因子を $E_0=q^{-1}(Z)$，tautological sectionを $\tau$ と書く。
線束の局所フレームでは $\tau=h\ell$ であり，
方向は $D=hD_0$ または $P=hP_0$ と書ける。
方向空間上の標準核束を $K$ とし，
\[
            \mathcal P_{e_d}=\operatorname{Sym}^{e_d}K
\]
と置く。これは $Y$ 上の係数束であると同時に，
ambient代数の中央な内部次数 $e_d$ の部分束を与える。
以下では $A,\Omega,\lambda$ の $Y$ への引戻しに同じ記号を用いる。

\begin{lemma}[規格化したambient中心の普遍的な核]
\label{lem:normalized-ambient-center-near-zero}
$V$ を $Z$ 全体を保った不変開集合へさらに縮小して，次を成立させられる。
生成元との交換子を $\tau$ で割った写像を
\[
 \mathcal D:A_{e_d}\longrightarrow
   \bigoplus_w\operatorname{Hom}
       (\mathcal V_w,A_{e_d+w})\otimes L^*
\]
と書くと，任意の可換基底変更 $T\to Y$ に対して
\begin{align}
 \ker(\mathcal D_T,\lambda_T)&=(\mathcal P_{e_d})_T,
       \label{eq:normalized-ambient-center-kernel}\\
 \ker\mathcal D_T
       &=(\mathcal P_{e_d})_T\oplus\mathcal O_T\Omega_T
       \label{eq:normalized-ambient-center-splitting}
\end{align}
が成立する。
両右辺は $(A_{e_d})_T$ の部分束である。
$\tau$ が可逆な基底では，第二式の左辺は通常のambient中心の次数 $e_d$ 部分に等しい。
零切断では通常の交換子の核はより大きく，第二式は除した交換子についての主張である。
\end{lemma}
\begin{proof}
$Y$ は滑らかな係数空間と方向空間上の線束の全空間の開部分であり，
例外因子は有効Cartier因子である。
$A_m$ は $Y$ 上有限局所自由なので，局所方程式 $h$ は各次数部分の非零因子である。
交換子は $h=0$ で消えるため，一意的に $h$ で割れる。
フレーム変更の下でこれらは $L^*$ の捩れを付けて貼り合い，$\mathcal D$ を与える。
この写像をまず普遍基底上で定めてから引き戻すので，
基底変更後に $h$ が零因子となる場合にも $\mathcal D_T$ は定義される。

$K$ の元は関係式により厳密に中央である。
完全対称化の同型から $\mathcal P_{e_d}\subset A_{e_d}$ は部分束であり，
$\mathcal D$ と $\lambda$ はその上で零である。
従って商束上の有限行列
\[
 j:\ A_{e_d}/\mathcal P_{e_d}\longrightarrow
   \left(\bigoplus_w\operatorname{Hom}
       (\mathcal V_w,A_{e_d+w})\otimes L^*\right)\oplus\mathcal O_Y,
 \qquad j=(\overline{\mathcal D},\overline\lambda)
\]
を得る。

例外因子の全幾何学的点で $\mathcal D$ は古典Poisson交換子である。
その内部次数 $e_d$ の中心は
\[
            \operatorname{Sym}^{e_d}K\oplus\mathcal O F,
 \qquad \lambda(F)=1/2
\]
である。これには，滑らかな可換中心商 $B_f=A_f/(F)$ の
Poisson中心が $\operatorname{Sym}K$ であることを用いればよい
（補題\ref{lem:classical-casimir-center}および
定理\ref{thm:low-center-zero-neighborhood}の全方向での確認）。
実際，Poisson中央な $a\in(A_f)_{e_d}$ の像は
$p\in\operatorname{Sym}^{e_d}K$ に属し，従って
$a-p=cF$，$c\in\mathbb C$ となる。
$\lambda(F)=1/2$ により規格化係数は残る $F$ 方向を検出する。
従って $j$ は例外因子の全点でファイバー上単射である。

有限行列の最大小行列式によって，$j$ が局所分裂単射となる開集合
$Y^{\mathrm{cen}}\subset Y$ は $E_0$ 全体を含む。
その閉補集合のproperな像を $V$ から除くと，零切断を失わずに
$q^{-1}(V)\subset Y^{\mathrm{cen}}$ とできる。
全ての写像が座標変換に自然なので，この開集合は不変に取れる。
以後この開集合へ制限する。

分裂単射は任意の基底変更後も単射であるから，
$\mathcal D_T a=0$ かつ $\lambda_T(a)=0$ ならば $a\in(\mathcal P_{e_d})_T$ である。
これが第一式を与える。
中心元 $\Omega$ は普遍基底上で $\mathcal D\Omega=0$ を満たす。
実際，通常の交換子は零であり，普遍基底では $h$ を消去できる。
この等式は任意基底変更後も成り立つ。
$\lambda(\Omega)=1/2$ なので，任意の $a\in\ker\mathcal D_T$ について
$a-2\lambda_T(a)\Omega_T\in(\mathcal P_{e_d})_T$ となり，第二式を得る。
部分束であることは，$\mathcal P_{e_d}\subset A_{e_d}$ が部分束であり，
商束上の $\Omega$ が $\lambda$ によって分裂することから従う。
\end{proof}

\begin{proposition}[Casimir係数の可逆なJacobianを持つ零切断近傍]
\label{prop:casimir-jacobian-zero-neighborhood}
前補題の開集合をさらに $Z$ 全体を保って不変に縮小すると，
次の係数方向微分は任意の方向解消点で同型となる：
\[
          J_\Omega:\mathcal P_{e_d}\xrightarrow{\ \sim\ }\mathcal P_{e_d}.
\]
この写像は，次数1では $A_4,D$ を固定して $A_6$ に
$c\in\operatorname{Sym}^6K$ を加える変化，次数2では $P$ を固定して
$U$ に $c\in\operatorname{Sym}^4K$ を加える変化に対する $\Omega$ の微分である。
その例外因子上の制限は恒等写像であり，同型性は任意の非被約基底変更後も保たれる。
\end{proposition}
\begin{proof}
まずこれらの係数変化の下でambient代数が文字通り不変であることを確認する。
次数1の表示関係式では $A_6$ は $DA_6$ としてだけ現れる。
$c\in\operatorname{Sym}^6K$ ならば $D_0c=0$ なので，$D(A_6+c)=DA_6$ である。
次数2では $U$ は $\sum_jp_{ij}\partial_{x_j}U$ としてだけ現れ，
$c\in\operatorname{Sym}^4K$ に対してこの方向微分は零である。
従って両次数とも，生成元を恒等に送るambient代数の標準同一視がある。
これは多項式の恒等式であり，$h=0$ や非被約基底でも成立する。
完全対称化写像と $\lambda$ もこの同一視で一致する。

係数変化を定義域に留める条件は開条件である。
特に平方零の変化 $c\varepsilon$ は同じ位相的点にあるので，
全ての既存の開条件を自動的に保つ。
上の標準同一視を用い
\[
 \Omega(f+\varepsilon c)
       =\Omega(f)+\varepsilon\,\partial_c\Omega,
       \qquad\varepsilon^2=0
\]
と定義すると，$\partial_c\Omega$ は $c$ に関して線型であり，
束の射 $\mathcal P_{e_d}\to A_{e_d}$ を与える。

除した交換子 $\mathcal D$ もこの係数変化に沿って一定である。
これはambientの同一視と交換子の一致を，普遍係数変化の空間上で $h$ により割って確認できる。
その普遍空間は $Y$ 上のベクトル束の開部分なので $h$ はなお非零因子であり，
その後に任意基底へ引き戻せる。
従って $\mathcal D\Omega=0$ と $\lambda(\Omega)=1/2$ を微分すると
\[
       \mathcal D(\partial_c\Omega)=0,
       \qquad\lambda(\partial_c\Omega)=0
\]
となる。前補題から $\partial_c\Omega\in\mathcal P_{e_d}$ であり，
これが内在的なendomorphism $J_\Omega$ を定める。
基底変更後に $h$ を改めて割っているのではない。

例外因子上では選択した中心元は
\[
 \Omega=F=
 \begin{cases}
 z^2/2+y^3/3+A_4y+A_6,&d=1,\\
 z^2/2+U,&d=2.
 \end{cases}
\]
従って上の係数方向微分は $c\mapsto c$ であり，
$J_\Omega|_{E_0}=\operatorname{Id}$ である。
$\det J_\Omega\ne0$ は $E_0$ 全体を含む開条件である。
再びその閉補集合のproperな像を $V$ から除くことで，
零切断を保ち，その全逆像で $J_\Omega$ を可逆にできる。
中心元を含む全構成は $\operatorname{GL}$ 同変であり，
このendomorphismの行列式条件も同変なので，開集合は不変に取れる。
得た束の同型を任意の基底へ引き戻せば，最後の主張も従う。
\end{proof}

\begin{lemma}[小拡大上の規格化中心差の一意な補正]
\label{lem:casimir-center-small-extension-correction}
前命題の開集合の非零方向部分で作業する。
$R'\twoheadrightarrow R$ を局所Artin $\mathbb C$-代数の平方零拡大，
$I$ をその核とし，$a':\Spec R'\to V^\times$ を一つの係数持上げとする。
そのambientを $A_{a'}$ とする。
同じambientの中心に属する内部次数 $e_d$ の元 $\Theta$ が
\[
          \lambda(\Theta)=1/2,
          \qquad \Theta\bmod I=\Omega_{a'\bmod I}
\]
を満たすならば，一意的な $\xi\in I(\mathcal P_{e_d})_{a'}$ があり，
Casimir係数だけを $\xi$ だけ動かした持上げ $a'+\xi$ に対して
\[
                    \Omega_{a'+\xi}=\Theta
\]
となる。両ambientは生成元を恒等とする標準同一視で比較する。
\end{lemma}
\begin{proof}
非零方向上では $q$ は同型であり，例外因子の局所方程式 $h$ は単元である。
従って通常の中心と除した交換子の核は一致する。
差を $\Delta=\Theta-\Omega_{a'}$ と置くと，
任意基底変更に両立する式\eqref{eq:normalized-ambient-center-kernel}より
\[
          \Delta\in(\mathcal P_{e_d})_{a'}.
\]
また還元が零なので $\Delta\in I(A_{e_d})_{a'}$ である。
$\mathcal P_{e_d}\subset A_{e_d}$ は部分束であり，商も有限局所自由であるため，
\[
 (\mathcal P_{e_d})_{a'}\cap I(A_{e_d})_{a'}
             =I(\mathcal P_{e_d})_{a'}.
\]
実際，局所Artin環上でこの部分束を直和因子として分裂すれば明らかである。
従って $\Delta$ は単に中心的な $I$-係数元というだけでなく，
厳密に $I(\mathcal P_{e_d})_{a'}$ に属する。

$J_{\Omega,a'}$ は同型なので
\[
               \xi=J_{\Omega,a'}^{-1}(\Delta)
                   \in I(\mathcal P_{e_d})_{a'}
\]
と取れる。
$I^2=0$ であるから，正則関数のTaylor公式は一次で厳密に止まり，
ambientの標準同一視の下で
\[
       \Omega_{a'+\xi}
          =\Omega_{a'}+J_{\Omega,a'}(\xi)
          =\Theta
\]
となる。
係数変化は平方零であり，既存の全ての開条件を保つ。
また方向と $K$ は変わらず，ambientの表示関係式も変わらない。
同じ等式を満たす別の補正との差は $J_{\Omega,a'}$ の核に入るので零である。
\end{proof}

\begin{remark}[この補正が与える範囲]
この議論は，有限中心元がCasimir係数に線型であることを示してはいない。
異なる有限係数点が同じ規格化中心元を与える離散的な分枝も排除しない。
示したのは，全可換零切断を保った一つの有限開近傍で，
この係数方向の方程式が\'{e}taleであり，平方零持上げを一意に補正できることである。
従って全同値群oidのsource持上げには利用できるが，
それだけから座標作用群oidの全射性・同型性を結論することはしない。
\end{remark}

\subsection{非零方向における次数付き同型と中央係数の差}

以下では全可換零切断を保った不変開集合 $W\subset V$ を取り，
系\ref{cor:zero-neighborhood-full-graded-recognition}の通常次数付き認識と，
補題\ref{lem:normalized-ambient-center-near-zero}，
命題\ref{prop:casimir-jacobian-zero-neighborhood}を同時に満たすようにする。
前項の有限個の完全対称化同型の条件も含める。
特にambientの各次数 $m\le e_d$ は，完全対称化によって
固定された可換重み付き多項式の次数 $m$ の部分と同一視できる。
$m<e_d$ では $B_m=A_m$ である。
係数環は局所Artin $\mathbb C$-代数 $\Lambda$ とし，
非零方向とは剰余体で方向が非零，従ってある方向座標が $\Lambda$ の単元であることを意味する。

\begin{lemma}[非零方向上の重み付き線型化]
\label{lem:nonzero-graded-isom-weighted-linear}
非零方向の二点 $a,b:\operatorname{Spec}\Lambda\to W^\times$ と
次数付き同型 $\phi:B_a\xrightarrow{\sim}B_b$ を取る。
その生成元上の写像は次の形である。
\begin{enumerate}
\item 次数2では，$M\in\operatorname{GL}_{n+1}(\Lambda)$ と
      $c\in\Lambda^\times$ によって
      $\phi(x)=Mx'$，$\phi(z)=cz'$ と書ける。
\item 次数1では，$M\in\operatorname{GL}_n(\Lambda)$ と
      $\beta,\gamma\in\Lambda^\times$ によって
      \[
       \phi(x)=Mx',\qquad \phi(y)=\beta y',\qquad
       \phi(z)=\gamma z',\qquad \gamma^2=\beta^3
      \]
      と書ける。特に $t=\gamma/\beta$ とすれば
      $(\beta,\gamma)=(t^2,t^3)$ である。
\end{enumerate}
これらは有効座標変換群 $G_d$ の元 $g$ を定める。
次数スカラーで $\phi$ を変更しても $g$ は変わらない。
標準の座標同型によってsourceを $ga$ へ移し，生成元上の写像を恒等に直すと，
targetとの差は次の中央係数だけとなる：
\[
 \begin{array}{c|c|c}
 d&\text{一致するパラメータ}&\text{残る係数差}\\\hline
 1&D,\ A_4&A_{6,b}-A_{6,ga}\in\operatorname{Sym}^6K,\\
 2&P&U_b-U_{ga}\in\operatorname{Sym}^4K.
 \end{array}
\]
両ambientは生成元を恒等とする同一視を持ち，その下で
\begin{equation}\label{eq:central-coefficient-arrow-equation}
                       \Omega_b=\Omega_{ga}
\end{equation}
となる。逆に，このパラメータの一致，中央係数差，中心元の等式は，
その重み付き線型座標写像から通常の次数付き同型を与える。
\end{lemma}
\begin{proof}
\emph{次数2。}
次数1では生成元 $x_i$ だけがあるため，$\phi(x)=Mx'$，$M$ は可逆である。
次数2の完全対称化分解から
\[
                \phi(z)=cz'+\operatorname{Sym}(q_2(x'))
\]
と書ける。
$\phi$ は $\operatorname{Sym}^2\langle x_i\rangle$ の部分加群を
$\operatorname{Sym}^2M$ によって可逆に保つ。
従ってその商 $\Lambda z$ に誘導する写像も可逆であり，$c\in\Lambda^\times$ である。
これは $z$ が非零方向で代数的に分解可能であることとは別の，線型部分空間の議論である。
関係式 $[x_i,x_j]=p_{ij}z$ に $\phi$ を適用すると
\[
       (MP_bM^{\mathsf t})_{ij}z'
            =p_{a,ij}\bigl(cz'+\operatorname{Sym}(q_2)\bigr)
\]
となる。
ある $p_{a,ij}$ は単元であり，次数2の直和分解から $q_2=0$ が従う。

\emph{次数1。}
同様に $\phi(x)=Mx'$ であり，次数2と3の完全対称化分解を使うと
\[
 \phi(y)=\beta y'+q_2(x'),\qquad
 \phi(z)=\gamma z'+\operatorname{Sym}(\ell_1(x')y')+q_3(x')
\]
と書ける。
純粋な $x$ の次数2部分を $\phi$ は可逆に保つから，その商により $\beta$ は単元である。
次数3の部分
$\operatorname{Sym}^3\langle x_i\rangle\oplus
\operatorname{Sym}(\langle x_i\rangle y)$ も
$M$ と $\beta$ による三角な可逆写像で保たれる。
その商は $\Lambda z$ なので，$\gamma$ も単元である。

関係式 $[x_i,y]=D_i z$ から
\[
 \beta(MD_b)_i z'
   =D_{a,i}\bigl(\gamma z'+\operatorname{Sym}(\ell_1y')+q_3\bigr)
\]
を得る。ある $D_{a,i}$ は単元なので
$\ell_1=q_3=0$，$\beta MD_b=\gamma D_a$ となる。
次に $[z,x_i]=D_i(y^2+A_4)$ を比較し，その単元 $D_{a,i}$ を消去すると
\begin{equation}\label{eq:weighted-linear-degree-four-comparison}
 \frac{\gamma^2}{\beta}\bigl((y')^2+A_{4,b}(x')\bigr)
   =(\beta y'+q_2(x'))^2+A_{4,a}(Mx')
\end{equation}
となる。

量子並べ替えの補正を省略せずに $q_2=0$ を確認する。
targetの導分を $D_b$ とすると，任意の二次式 $q=q_2(x')$ に対し
\[
 y'q+qy'
  =2\operatorname{Sym}(qy')
       -\frac{D_b^2q}{6}\bigl((y')^2+A_{4,b}\bigr).
\]
これは $q=x_i'x_j'$ について三つの置換位置を平均し，
$[x_i',y']=D_{b,i}z'$ と
$[z',x_i']=D_{b,i}((y')^2+A_{4,b})$ を二回使えば得られ，一般の場合は線型性による。
従って\eqref{eq:weighted-linear-degree-four-comparison}の
$\operatorname{Sym}(\operatorname{Sym}^2\langle x_i'\rangle\, y')$ 成分は
右辺の $2\beta q_2$ だけである。
完全対称化は次数4で同型，$2\beta$ は単元なので $q_2=0$ となる。
その後に $(y')^2$ 成分を比較して $\gamma^2=\beta^3$ を得る。

\emph{残るCasimir係数。}
得た写像はいずれも正規形に許された重み付き線型座標変換である。
第\ref{sec:definitions}節の作用の規約に従い，それが誘導する同型を
$B_a\to B_{ga}$ と書く。
その逆を合成すれば，全生成元を恒等に送る同型の問題となる。
次数スカラーは $G_d$ を定義する中央部分群に一致するから，$g$ はその商に依存しない。

次数2ではまず $P$ が一致し，残る関係式から
\[
     \sum_jp_{ij}\partial_{x_j}(U_b-U_{ga})=0
\]
を得る。
この等式はまず $\operatorname{Sym}$ による像についての等式だが，
関係式の次数3は中央関係式の次数4より小さく，
ambientの次数3の完全対称化は単射なので，多項式として成立する。
局所的に $K$ の基底と活性座標 $u,v$ を選べば，
$P(u,v)$ は単元であり，上式は $\partial_u(U_b-U_{ga})=
\partial_v(U_b-U_{ga})=0$ と同値である。
係数環は $\mathbb C$-代数なので，非被約でも
$U_b-U_{ga}\in\operatorname{Sym}^4K$ となる。

次数1では順に $D,A_4$ が一致する。
最後の関係式 $[y,z]=\operatorname{Sym}((DA_4)y)+DA_6$ を比較すると
$D(A_{6,b}-A_{6,ga})=0$ となる。
次数5の完全対称化の単射性を用いて多項式の等式に戻し，
活性座標に沿う導分の核を取れば
$A_{6,b}-A_{6,ga}\in\operatorname{Sym}^6K$ を得る。
いずれもこの係数差はambientの全表示関係式を変えない。

ambient関係式の次数は全て $e_d$ より小さい。
従って生成元の恒等写像は共通ambientの恒等写像へ持ち上がり，
quotientの同型は二つの中心元が生成するイデアルの一致を意味する。
内部次数 $e_d$ のそのイデアルはそれぞれ $\Lambda\Omega_b$ と
$\Lambda\Omega_{ga}$ である。
従って中心元は単元倍で一致するが，両者の $z^2$ 係数が $1/2$ なのでその単元は1である。
これが\eqref{eq:central-coefficient-arrow-equation}である。
逆向きは共通ambientの同じ中心商を取ることから従う。
\end{proof}

\subsection{全ての矢を保った零切断近傍でのsourceの持上げ}

\begin{theorem}[可換零切断近傍の全同値群oidの滑らかさ]
\label{thm:full-source-smooth-zero-neighborhood}
滑らかな可換零切断全体を含む不変な有限型開集合 $W_d\subset V_d$ を選び，
次を同時に満たすようにできる。
\begin{enumerate}
\item 通常のgraded同型による全認識
      $\mathcal H_{W_d}\simeq\mathcal E_{W_d}$ が成立する。
\item 全ての矢を保った群oid
      $\mathcal E_{W_d}\rightrightarrows W_d$ のsource・targetは滑らかであり，
      相対次元は $g_1=n^2$，$g_2=(n+1)^2$ である。
\item 非零方向を結ぶ矢上で，前補題の重み付き線型成分を記録する射
      \[
       \Psi:\mathcal E_{W_d^\times}\longrightarrow G_d\times W_d^\times,
       \qquad \alpha\longmapsto(g(\alpha),s(\alpha))
      \]
      は\'{e}taleである。
\end{enumerate}
\end{theorem}
\begin{proof}
系\ref{cor:zero-neighborhood-full-graded-recognition}，
補題\ref{lem:normalized-ambient-center-near-zero}，
命題\ref{prop:casimir-jacobian-zero-neighborhood}および低次数完全対称化の
有限個の開集合の共通部分を $W_d$ とする。
いずれも零切断全体を保つため，その共通部分も同じ性質を持つ。
(1)は系\ref{cor:zero-neighborhood-full-graded-recognition}による。
前補題の低次数直和比較はArtin性を使わず，任意の基底の方向座標が局所的に可逆なchartで成立する。
従って普遍的な通常次数付きIsom上でも非線型成分は恒等的に零であり，
低次数係数を読む操作は全基底に自然な代数的操作である。
これは次数スカラーを法として $G_d$ に値を取る。
従って (3) の射 $\Psi$ が定義される。
その形式的\'{e}tale性を，閉点を特殊点とする局所Artin環の小拡大で確認する。

$\Lambda'\twoheadrightarrow\Lambda$ を小拡大，$I$ をその核とする。
非零方向の矢 $\alpha:a\to b$ と，
$\Psi(\alpha)=(g,a)$ の持上げ $(g',a')\in(G_d\times W_d^\times)(\Lambda')$ を与える。
$a_0'=g'a'$ と置く。
座標群の作用と開集合の不変性により，$a_0'$ も同じ有限族の定義域に入る。
前補題から $b-ga$ は共通方向のCasimir係数
$\zeta\in\operatorname{Sym}^{e_d}K_\Lambda$ だけである。
この束は $\Lambda'$ 上自由なので，$\zeta$ を
$\zeta_0'\in\operatorname{Sym}^{e_d}K_{\Lambda'}$ へ持ち上げ，
\[
                   b_0'=a_0'+\zeta_0'
\]
とする。
この点は $b$ へ還元するから，全ての開条件を満たす。
また $A_{a_0'}=A_{b_0'}$ は生成元を恒等とする共通ambientである。

このambient内で $\Theta=\Omega_{a_0'}$ と置く。
元の矢に対する\eqref{eq:central-coefficient-arrow-equation}より
$\Theta\bmod I=\Omega_b$ であり，$\lambda(\Theta)=1/2$ である。
補題\ref{lem:casimir-center-small-extension-correction}を実際のtargetの持上げ $b_0'$ に適用すると，
一意的な $\xi\in I\operatorname{Sym}^{e_d}K_{\Lambda'}$ があり，
\[
              b'=b_0'+\xi,
              \qquad\Omega_{b'}=\Theta
\]
となる。
従って座標同型 $B_{a'}\simeq B_{a_0'}$ と共通ambientの恒等を合成して，
$\alpha$ を持ち上げる矢 $a'\to b'$ を得る。

同じ $(g',a')$ を持つ別の持上げがあれば，前補題によりそのtargetとの差は
$I\operatorname{Sym}^{e_d}K$ に限られる。
中心差の補正の一意性からtargetは $b'$ と一致する。
その生成元上の写像は $g'$ で決まり，残る次数スカラーは既に割っているので，
矢の類も一致する。
従って $\Psi$ は閉点で形式的\'{e}taleである。
有限型射の閉点Artin持上げ判定では，存在が滑らかさを与え，一意性が相対接空間の消滅を与える。
よって $\Psi$ は全閉点で\'{e}taleであり，非\'{e}tale locusが閉であるため全点で\'{e}taleである
\cite[Tags 02HY, 02HM]{StacksProject}。

$s=\operatorname{pr}_{W_d^\times}\circ\Psi$ なので，
非零方向を結ぶ全矢でsourceは滑らか，相対次元は $\dim G_d=g_d$ である。
零切断に接する矢は補題\ref{lem:commutative-zero-equivalence-saturated}により両端が零切断にあり，
定理\ref{thm:equivalence-source-formal-smooth}が同じ滑らかさと相対次元を与える。
これで全複素閉矢においてsourceが滑らかであり，
非smooth locusが閉であることから全ての点で滑らかである。
逆射がsourceとtargetを交換するため，targetにも同じ結論が従う。
\end{proof}

\begin{remark}[座標作用への全射性は必要ない]
\label{rem:central-coefficient-branches}
ここで証明したのは $\Psi$ が\'{e}taleであることであり，
任意の矢のtargetが厳密に $ga$ と一致することではない。
選んだ有限中心元は中央係数に非線型に依存してよく，
同じ規格化中心を与える離散的な別の係数分枝が残り得る。
その分枝はsourceの小拡大持上げを妨げない。
持上げで使ったのはtarget側のCasimir Jacobianの可逆性であり，
中央係数の大域的単射性や全係数域での一意性ではない。
\end{remark}

\subsection{全可換零切断とその形式変形を含むDM部分スタック}

\begin{theorem}[全可換零切断を含む有限型DM部分スタック]
\label{thm:zero-neighborhood-full-dm}
前定理の全同値群oidの商
\[
             \mathfrak U_{n,d}
                      =[W_d/\mathcal E_{W_d}]
\]
は，有限型で整なDeligne--Mumfordスタックである。
その対角はaffine，有限表示，quasi-finiteかつunramifiedであり，
$W_d\to\mathfrak U_{n,d}$ は相対次元 $g_d$ の滑らかな全射である。
次元は
\[
 \dim\mathfrak U_{n,1}=b_{n,1}+n,
 \qquad
 \dim\mathfrak U_{n,2}=b_{n,2}+2n-1
\]
である。
次数1，および次数2の $n=2$ ではこのスタックは $\mathbb C$ 上滑らかである。
次数2の $n\ge3$ では，滑らかでない点の集合は可換零切断の像と一致し，
その被約閉部分スタックは $[Z/\mathcal E_Z]$ である。
この主張は非smooth locusの被約構造についてのものであり，Jacobianイデアルによる構造との一致は主張しない。

任意の非空対象開 $W'\subset W_d$ に対して，全群oidを制限した商
$[W'/\mathcal E_{W'}]$ は $\mathfrak U_{n,d}$ の非空開部分スタックとなる。
従って後述の系\ref{thm:generic-full-dm-stack}の一般非零部分 $\mathfrak N_{n,d}$ は，
$\mathfrak U_{n,d}$ の非空開部分スタックとして取れる。

許された局所Noether基底のsite上では，この商は
$\mathfrak M^{\mathcal O}_{n,d}$ のfullな部分スタックであり，
対象はfppf局所的に $W_d$ 上の構成族から来る族，射はその間の全構造対象付きAbel同値である。
全ての滑らかな可換模型と，その全framed Abel Artin変形を含む。
この包含の開性は後の定理\ref{thm:zero-neighborhood-open-full-moduli}で証明する。
さらに，零切断全体を保つ縮小後の有限対角は
定理\ref{thm:zero-neighborhood-finite-diagonal}，
可換模型への表現可能な射は命題\ref{prop:representable-commutative-shadow}による。
全非可換対象の被覆は主張しない。
\end{theorem}
\begin{proof}
前定理と滑らかな群oidの商定理\cite[Tag 04TK]{StacksProject}により代数スタックを得る。
系\ref{cor:finite-equivalence-unramified}を表示へ引き戻し，忠実平坦降下すると，
対角の性質を得る。特に商はDeligne--Mumfordである
\cite[Tag 06N3]{StacksProject}。
有限型性は $W_d$ の有限型性と準コンパクト性から従う。

$W_d$ はaffine係数空間の開集合と既約方向錐の積の非空開集合であり，整である。
既約性は表示の全射性から商へ移り，被約性も滑らかな表示上で検査できるため商は整である。
係数空間の次元は $b_{n,d}+g_d$，方向錐の次元は次数1で $n$，次数2で $2n-1$ である。
表示の相対次元 $g_d$ を引けば表示した次元を得る。

次数1の方向空間はaffine空間である。
次数2で $n=2$ のPl\"ucker錐もaffine空間であり，$n\ge3$ では頂点だけが特異である。
従って $W_d$ が滑らかでない点は，後者の場合に限ってちょうど零切断である。
基底体上の滑らかさは滑らかな表示上で検査できるので，商についての主張が従う。
系\ref{cor:zero-section-scheme-saturated}により零切断のイデアルは群oid上で一致し，
$Z$ は閉部分スタック $[Z/\mathcal E_Z]$ へ降下する。
その表示 $Z$ は滑らかであるためこの閉部分スタックは被約であり，
述べた非smooth locusの被約構造との一致を得る。

sourceとtargetは滑らかであるから開写像であり，任意の対象開 $W'$ の飽和
$s(t^{-1}(W'))$ は開集合である。
fullな群oid制限の商はこの飽和開集合の商と同じであるから，
$[W'/\mathcal E_{W'}]\to\mathfrak U_{n,d}$ は開埋込みとなる。

fullな部分問題との同定には，定理\ref{thm:finite-anticanonical-isom}の
全同値層とtailsのfppf降下を用いる。
局所的な対象は定義によって $W_d$ の点となり，その間の全射が
$\mathcal E_{W_d}$ であるから，商のスタック化はちょうど述べたfullな部分スタックである。
零切断全体が $W_d$ に含まれ，全Artin変形は
定理\ref{thm:simultaneous-central-all-artin}のsliceのArtin点から来る。
その特殊点が $W_d$ に入れば，Artin点全体もこの開集合を通る。
従って最後の包含も得られる。
\end{proof}

\subsection{可換点における有限型の\'{e}tale slice}

\begin{corollary}[有限自己同型群を保つ幾何学的slice]
\label{cor:finite-etale-slice-near-classical}
定理\ref{thm:zero-neighborhood-full-dm}の有限型DMスタックを
\[
 \mathfrak U_{n,d}=[W_d/\mathcal E_{W_d}]
\]
とし，滑らかな可換点 $x=(f,0)\in W_d(\mathbb C)$ の多様体を $X$ とする。
有限群 $\Gamma=\operatorname{Aut}(X)$ を，正規形の有効座標変換群 $G_d$ における
$f$ の安定化群と同一視する。
全係数空間の軌道接空間に対し，$\Gamma$ 不変な補空間
\[
 T_f\mathcal F_d=(\operatorname{Lie}G_d)\cdot f\oplus N,
 \qquad \dim N=b_{n,d}
\]
を選べる。
さらに $(f,0)$ を含む $\Gamma$ 不変なaffine開近傍
\[
 T\subset (f+N)\times C_{n,d}
\]
を取り，次を満たすようにできる。
\begin{enumerate}
\item $T\subset W_d$ であり，同じ有限中心商族の制限は
      representableな\'{e}tale射
      \[
       \eta_X:[T/\Gamma]\longrightarrow\mathfrak U_{n,d}
      \]
      を定める。
\item $(f,0)$ の安定化群から像の安定化群への射は
      $\Gamma\xrightarrow{\sim}\operatorname{Aut}(X)$ である。
      従って残余gerbe $B\Gamma$ 上で同型を与える。
\item $\eta_X$ の像は $[X]$ を含む
      $\mathfrak U_{n,d}$ の開部分スタックであり，
      $[T/\Gamma]$ はその有限型\'{e}tale被覆となる。
      その $B\Gamma$ に沿う形式完備化は，既存の
      $[\operatorname{Spf}R_X/\Gamma]$ と両立する。
\end{enumerate}
ここでrepresentableとは代数空間によって表現可能という意味である。
この系は構成済みのfullな部分スタックの内部での幾何学的局所表示であり，
全 $\mathfrak M^{\mathcal O}_{n,d}$ の中での開性には
後の定理\ref{thm:zero-neighborhood-open-full-moduli}を用いる。
\end{corollary}
\begin{proof}
\emph{有限群に不変な横断slice。}
正規形の商表示と可換点での全 $\mathcal O$ 保存同値の認識により，
$\Gamma$ は $G_d$ における $f$ の有限な安定化群である。
$H^0(X,T_X)=0$ なので軌道接写像は単射である。
軌道接空間は $\Gamma$ 不変であり，標数0の有限群の平均化によって
その不変な補空間 $N$ を選べる。
これは命題\ref{prop:equivariant-formal-product}と同じ補空間の選択である。
$G_d$ の係数空間への作用は線型で $\Gamma$ は $f$ を固定するので，
$f+N$ は $\Gamma$ 不変なaffine部分空間である。

滑らかな可換係数開集合を $U_d$ とし，
\[
 S=(f+N)\cap U_d,\qquad
 T_0=(S\times C_{n,d})\cap W_d
\]
と置く。
これは $\Gamma$ 不変で，$x=(f,0)$ を含む。
座標作用の射\eqref{eq:coordinate-action-to-equivalences}を $\Gamma$ へ制限すると，
同じ有限族の写像
\[
 p:T_0\longrightarrow\mathfrak U_{n,d}
\]
には積と両立する $\Gamma$ 同変構造がある。
従って $[T_0/\Gamma]$ からの射 $\overline p$ を得る。
次数スカラーが与える自然同型のcoherenceは，既存の有効群作用の構成に含まれている。

\emph{sliceからの射の\'{e}tale性。}
滑らかな表示 $W_d\to\mathfrak U_{n,d}$ で基底変更すると，
\[
 P=T_0\times_{\mathfrak U_{n,d}}W_d
   \simeq\mathcal E_{W_d}\times_{s,W_d}T_0
\]
を得る。これはスキームであり，$P\to W_d$ はtargetから誘導される。
単位矢に対応する点を $e_x\in P$ とする。
補題\ref{lem:full-coefficient-formal-slice}と
定理\ref{thm:equivalence-source-formal-smooth}の，同じsliceに適合した同型を使うと
\[
 \widehat{\mathcal O}_{W_d,x}=R_X[[u_1,\ldots,u_{g_d}]],
 \qquad
 \widehat{\mathcal O}_{T_0,x}=R_X
\]
であり，sliceの埋込みは $u=0$ に対応する。
また矢の完備局所環は，sourceとtargetの写像に両立して
\[
 \widehat{\mathcal O}_{\mathcal E,e_x}
       =R_X[[u_1,\ldots,u_{g_d},v_1,\ldots,v_{g_d}]]
\]
である。
従って $P$ の単位点における完成は
\[
 \widehat{\mathcal O}_{P,e_x}
   =R_X[[u,v]]\widehat\otimes_{R_X[[u]]}R_X
   \simeq R_X[[v]],
\]
となり，$P\to W_d$ はtargetの完成
$\widehat{\mathcal O}_{W_d,x}=R_X[[v]]$ からの同型を誘導する。

この完成計算は，局所Artin環に対する関手としても確認できる。
任意のtargetのパラメータ持上げに対し，そのframed変形類は
$T_0$ の完備化の一意な点で表され，特殊ファイバーで恒等となる矢も一意である。
従って上の同型は抽象的な環同型ではなく，実際のtarget写像に沿う同型である。
特に，次数2のPl\"ucker錐が特異でもこの議論は有効である。

$P\to W_d$ は有限型複素スキームの射であるから，
完備局所環の判定\cite[Lemma~41.11.3, Tag~039M]{StacksProject}により
$e_x$ で\'{e}taleである。
representableな射の\'{e}tale性は滑らかな表示によって検査できるため，
$p$ は $x$ で\'{e}taleである
\cite[Section~100.3, Tag~04XB]{StacksProject}。
従って $x$ を含む開近傍 $T_1\subset T_0$ で $p$ は\'{e}taleとなる。

有限群による全ての移動の共通部分を取れば，$T_1$ を $\Gamma$ 不変にできる。
さらに $(f+N)\times C_{n,d}$ はaffineであり，
$T_1$ はその $\Gamma$ 固定点 $x$ を含む開集合である。
$x\in D(a)\subset T_1$ となる正則関数 $a$ を選び，
\[
                         a_\Gamma=\prod_{\gamma\in\Gamma}\gamma(a)
\]
と置けば，$T=D(a_\Gamma)$ は $x$ を含む $\Gamma$ 不変なaffine開集合で，
$T\subset T_1$ となる。

\emph{商からの射のrepresentability。}
$\eta_X:[T/\Gamma]\to\mathfrak U_{n,d}$ の
相対的な自己同型の核が自明であることを示す。
任意の基底上で，$\Gamma$-torsorをfppf局所的に自明化すれば，
$[T/\Gamma]$ の対象は $t:T'\to T$ として表され，
その自己同型は $t$ を固定する $\gamma\in\Gamma(T')$ である。
その像が $\mathcal O$-marked tails上で恒等なら，
補題\ref{lem:ordinary-graded-recognition-injective}の核の計算により，
対応する次数付き座標変換はfppf局所的に次数スカラーである。
次数スカラーは有効群 $G_d$ で恒等である。
$\Gamma\hookrightarrow G_d$ は単射なので $\gamma=1$ であり，
相対自己同型の核は任意基底で自明である。
従って $\eta_X$ は代数空間によって表現可能である
\cite[Lemma~101.6.2, Tag~04YY]{StacksProject}。

$T\to[T/\Gamma]$ は有限\'{e}taleな全射であり，
その後の合成は既に\'{e}taleと示した $p|_T$ である。
representableな射の\'{e}tale性の降下から，$\eta_X$ 自身が\'{e}taleとなる。
固定点 $x$ では始域の安定化群は $\Gamma$ 全体であり，
終域の安定化群も古典的認識により $\operatorname{Aut}(X)=\Gamma$ なので，
残余gerbe上の主張を得る。

\'{E}tale射は開射であるから，その像は $[X]$ を含む
$\mathfrak U_{n,d}$ の開部分スタックとなる。
その像へ制限すれば $\eta_X$ は全射な\'{e}tale射である。
また $T$ への縮小は $x$ における完成を変えず，
同じ有限族のclassifying mapは $\Gamma$ 同変である。
従って形式完備化は
定理\ref{thm:unframed-formal-quotient}の
$[\operatorname{Spf}R_X/\Gamma]$ と両立する。
\end{proof}

\begin{remark}
このsliceは，抽象的な形式変形環の代数化を別に仮定して得たものではない。
元の全係数有限族の線型横断面を方向錐と掛け，有限個の開条件を課した実際の族である。
全モジュライ内での開性には定理\ref{thm:zero-neighborhood-open-full-moduli}を用いる。
有限対角を得た後には，系\ref{cor:stabilizer-preserving-slice-generic-trivial}によって
sliceを全点で安定化群を保つように縮められ，
系\ref{cor:algebraic-coarse-local-model}によって粗空間上の\'{e}tale局所模型も得られる。
これらは全非可換対象を被覆するという主張ではない。
また，この\'{e}tale射をZariski開埋込み，あるいは全構成族の大域的な有限群商とはしていない。
\end{remark}


\subsection{非零方向の滑らかな開部分スタック}

\begin{corollary}[全同値による一般非可換部分のDMモジュライ]
\label{thm:generic-full-dm-stack}
定理\ref{thm:zero-neighborhood-full-dm}の対象空間の非零方向部分で
任意の非空開集合 $W$ を取る。
全同値による商
\[
             \mathfrak N_{n,d}=[W/\mathcal E_W]
                 \longrightarrow\mathfrak U_{n,d}
\]
は非空開埋込みであり，$\mathfrak N_{n,d}$ は滑らかで整な有限型DMスタックである。
表示 $W\to\mathfrak N_{n,d}$ は相対次元 $g_d$ の滑らかな全射で，
対角はaffine，有限表示，quasi-finiteかつunramifiedである。
その次元は
\begin{align*}
 \dim\mathfrak N_{n,1}
   &=b_{n,1}+n=\binom{n+3}{4}+\binom{n+5}{6}-n^2+n,\\
 \dim\mathfrak N_{n,2}
   &=b_{n,2}+2n-1=\binom{n+4}{4}-(n+1)^2+2n-1.
\end{align*}
局所Noether fppf site上では，対象はfppf局所的に $W$ の構成族に同値な族であり，
射には全ての $\mathcal O$ 保存完全Abel同値を含める。
これらは通常の次数付き代数の同型を次数スカラーで割ったものとして認識できる。
\end{corollary}
\begin{proof}
定理\ref{thm:zero-neighborhood-full-dm}の対象開に関する結論を適用する。
全同値の制限であるため，同型層とfull essential imageも同じである。
非零方向空間は，次数1では $\mathbb A^n\setminus\{0\}$，
次数2では $\mathcal O_{\operatorname{Gr}(2,n+1)}(-1)$ の零切断を除いた全空間である。
従って $W$ は滑らかで，商も滑らかな表示上でこの性質を持つ。
次元と対角の性質は同定理から継承される。
通常の次数付き認識は系\ref{cor:zero-neighborhood-full-graded-recognition}による。
\end{proof}

\begin{remark}[部分問題の範囲]\label{rem:generic-dm-scope}
対象開 $W$ 自体が同値で飽和している必要はない。
$\mathfrak N_{n,d}$ は有限型DM部分問題 $\mathfrak U_{n,d}$ の開部分スタックである。
定理\ref{thm:zero-neighborhood-open-full-moduli}により，
零切断を保つ必要な縮小後の $\mathfrak U_{n,d}$ は全モジュライ
$\mathfrak M^{\mathcal O}_{n,d}$ 内で開となり，
その非零開部分 $\mathfrak N_{n,d}$ も同様である。
全非可換対象の被覆と，可換locusに接しない成分の有無は未解決である。
\end{remark}

\subsection{一般安定化群と粗モジュライ空間}

\begin{corollary}[有限\'{e}taleなinertiaと滑らかな粗空間]
\label{cor:generic-finite-inertia-coarse}
系\ref{thm:generic-full-dm-stack}の $W$ をさらに非空開集合へ縮小することで，
$\mathfrak N_{n,d}$ のinertiaを有限\'{e}taleにできる。
このとき $\mathfrak N_{n,d}$ は有限\'{e}tale群を帯とするgerbeとして
粗モジュライ代数空間
\[
       \pi:\mathfrak N_{n,d}\longrightarrow M_{n,d}
\]
を持ち，$M_{n,d}$ は $\mathbb C$ 上滑らか，整，有限型である。
特に正規であり，次元は系\ref{thm:generic-full-dm-stack}の次元式と同じである。
ここで帯は局所的に定まる有限\'{e}tale群であり，大域的に固定した群やgerbeの中立性は要求しない。
定理\ref{thm:zero-neighborhood-finite-diagonal}の縮小を先に行えば，
$M_{n,d}$ は分離的でもある。
さらに系\ref{cor:stabilizer-preserving-slice-generic-trivial}により，
この非空開部分を慣性自明な滑らかな代数空間として選べる。
\end{corollary}
\begin{proof}
まず $W=\Spec A$ を非空affine開集合に制限する。
その安定化群スキーム $I\to W$ はaffine，有限表示，quasi-finiteかつunramifiedである。
$I=\Spec B$，$K=\operatorname{Frac}(A)$ と書く。
有限型で零次元の $K$-代数 $B\otimes_AK$ は有限次元である。
$B$ の有限個の $A$-代数生成元 $b_i$ は，$K$ 上でmonic関係式を満たす。
その有限個の係数の分母を払い，さらに関係式の値として残る $A$-torsion元を殺す
非零元をinvertすれば，$B_f$ の各生成元は $A_f$ 上整となる。
従って $B_f$ は有限 $A_f$-加群である。
さらにgeneric flatnessにより非空開へ縮小し，$B_f$ を平坦にもできる。
有限表示かつunramifiedであることと合わせて，$I\to W$ は有限\'{e}taleとなる。

開集合への全群oid制限は，定理\ref{thm:full-source-smooth-zero-neighborhood}の
source・targetの滑らかさを保つ。
新しい商のinertiaをその滑らかな表示 $W$ に引き戻すと上の $I$ になるから，
有限\'{e}tale性は降下して $\mathcal I_{\mathfrak N_{n,d}}\to\mathfrak N_{n,d}$ に成立する。
Keel--Mori定理\cite[Tag 0DUT]{StacksProject}により粗モジュライ代数空間が存在し，
有限型性は\cite[Tag 0DUW]{StacksProject}から従う。

さらにflatかつ有限表示のinertiaを持つ代数スタックは，ある代数空間上のgerbeである
\cite[Tag 06QJ]{StacksProject}。
このgerbeの基底は幾何学的同型類と代数空間への射を分類するので，
粗空間の普遍性により $M_{n,d}$ と同一視される。
fppf局所的には $[T/G]\to T$ であり，$G\to T$ は有限\'{e}taleである
\cite[Tag 06QH]{StacksProject}。
$T\to[T/G]$ は有限\'{e}tale全射なので，$[T/G]\to T$ は相対次元零の滑らかな射である。
この性質を降下すると $\pi$ も相対次元零で滑らかかつ全射となる。
従って $\mathfrak N_{n,d}$ の $\mathbb C$ 上の滑らかさは $M_{n,d}$ に降下する。
また $\pi$ は幾何学的同型類を一つの点に移す全射であり，
既約性と次元も保たれる。
分離性は有限対角と
系\ref{cor:zero-neighborhood-separated-coarse}から従い，
慣性自明な開部分の存在は
系\ref{cor:stabilizer-preserving-slice-generic-trivial}で別に証明する。
\end{proof}

\begin{remark}[有限慣性と平坦な慣性の違い]
上のgeneric flatnessの議論は，一般のquasi-finiteな慣性を
有限\'{e}taleへ制限するための判定として有用である。
本族ではさらに定理\ref{thm:zero-neighborhood-finite-diagonal}と
系\ref{cor:zero-neighborhood-separated-coarse}により，
全可換零切断を含む近傍で有限慣性と分離的な粗空間を得る。
ただし可換点には被覆対合があり，
系\ref{cor:stabilizer-preserving-slice-generic-trivial}の一般点では安定化群が自明なので，
慣性の次数は零切断で跳ぶ。
従って零切断全体を保つ近傍で慣性を平坦または有限\'{e}taleとすることはできない。
一般開部分のgerbe表示をこの全近傍へ延長してはいない。
\end{remark}


\subsection{有限型群oidのlimit preservationと可換点での形式対象のeffectivity}

\begin{lemma}\label{lem:finite-groupoid-limit-preserving}
$\mathcal X_V=[V/\mathcal E]$ を，有限型の全同値群oidのfppf stackificationとする。
これは全 $\mathbb C$-スキーム上の抽象的な商stackとして定義できる。
その対角はアフィンかつ有限表示であり，$\mathcal X_V$ はlimit preservingである。
許されたNoether基底上では，構成族へfppf局所的に入る対象からなる
全 $\mathcal O$ 保存同値付きのfullな部分問題を表す。
一般の $V$ について $\mathcal X_V$ の代数性はこの補題の結論に含めない。
零切断近傍 $W_d$ への制限では，既に
定理\ref{thm:zero-neighborhood-full-dm}によって代数性を得ている。
\end{lemma}
\begin{proof}
二つの対象をfppf局所的に $V$ の点で表せば，その同型層は
$\mathcal E\to V\times V$ の基底変更になる。
アフィン射と有限表示性の忠実平坦降下により，対角の主張を得る。

$T=\operatorname{Spec}A$，$A=\varinjlim A_i$ をfiltered colimitとする。
$\mathcal X_V(T)$ の対象はfppf局所的に $V$ の点となる。
$T$ は準コンパクトでfppf射は開射なので，有限個のaffine開を選び，
一つの忠実平坦有限表示affine被覆 $T'\to T$ でその対象を表せる。
その降下データは $T'\to V$，$T'\times_TT'\to\mathcal E$ と
$T'\times_TT'\times_TT'$ 上の一つのcocycle等式である。
被覆，$V$，$\mathcal E$ は有限表示であるため，
これらの有限個の射と等式は十分大きい有限段階 $A_i$ へ降下する。
被覆の平坦性と全射性も，有限表示射の極限に関する通常の降下により
十分大きい有限段階で成立する。
従って対象は $\mathcal X_V(\operatorname{Spec}A_i)$ へ降下する。
射と射の等式は，既に確認した有限表示affine同型層によって有限段階へ降下する。
これで群oidの同値
\[
 \varinjlim_i\mathcal X_V(\operatorname{Spec}A_i)
     \simeq\mathcal X_V(\operatorname{Spec}A)
\]
を得る。
Noether基底上での解釈は，全同値の認識と既存のtailsのfppf降下による。
\end{proof}

\begin{proposition}[可換特殊ファイバーを持つ形式対象の実現]
\label{prop:commutative-formal-family-effective}
$X$ を滑らかな可換模型，$(R,\mathfrak m)$ を剰余体 $\mathbb C$ の
完備Noether局所 $\mathbb C$-代数とする。
$R_a=R/\mathfrak m^{a+1}$ 上の全framed Abel変形の両立系を与える。
この形式対象は，同時中心商の一つの固定した可換横断sliceの有限族を
$\operatorname{Spec}R$ へ引き戻した実際のnc del Pezzo族によって実現される。
さらに，構成族の二つの $R$ 上の模型の間では，
両立する全 $\mathcal O$ 保存形式同値は一意的に $R$ 上の同値へ実現される。
\end{proposition}
\begin{proof}
定理\ref{thm:simultaneous-central-all-artin}により，固定sliceの完備局所環
$C_X$ から各 $R_a$ への一意的な連続局所写像 $b_a$ を得る。
framed同値の一意性により，これらは $a$ に関して両立する。
$R=\varprojlim R_a$ なので，一意的な連続局所写像
$b:C_X\to R$ が存在する。
これを有限型sliceの局所環へ制限すると，
$\operatorname{Spec}R\to V$ を得る。
その特殊点は既存の全有限開条件を満たすので，必要な分母は $R$ の単元となる。
同じ有限族の引戻し $B_b$ は，既に証明した任意Noether基底変更の定理により
nc del Pezzo族を与え，その各 $R_a$ 上の還元は元の形式対象に同値である。
両立する同値はframed自己同型が自明であることから選べる。

構成族の二模型 $B_a,B_b$ の間の同値層は，
$\mathcal E$ の基底変更として $R$ 上affine有限表示である。
その座標環を $A$ とすれば
\[
 \operatorname{Hom}_{R\text{-alg}}(A,R)
    =\varprojlim_a\operatorname{Hom}_{R\text{-alg}}(A,R_a)
\]
である。従って両立する形式同値は一意的に実現される。
ここで一意性は標識を保つ自然同型を法とした意味であり，
その自然自己同型も既存の剛性により自明である。
\end{proof}

\begin{remark}[形式対象のeffectivityと全モジュライの比較]
前二結果は，構成族の像の有限型群oidについてのlimit preservationと，
可換特殊ファイバーを持つ形式対象のeffectivityを与える。
これらの証明だけでは任意の $\mathfrak M^{\mathcal O}_{n,d}(R)$ の対象や，
任意の二対象間の形式同値を扱っていない。
定理\ref{thm:zero-neighborhood-open-full-moduli}では，
任意のadmissible族との混合同値層と，構成族の全点での変形比較を別に用いて，
構成した近傍の全モジュライ内での開性を証明する。
全非可換対象の被覆や，可換locusに接しない部分の代数性はこの結論に含まれない。
一般にArtinの基準では，変形の比較，Rim--Schlessinger条件，effectivity，
versalityの開性などを分けて確認する必要がある
\cite[Tags 07T0, 07XP]{StacksProject}。
特に有限表示の射スキーム，有限な自己同値群，limit preservationだけを，
source・targetの平坦性や商stackの代数性の代用にはしない。
\end{remark}

\subsection{一般平坦性による補助判定}

ここでいう群oidの制限は，対象だけを開集合へ制限し，その間の全ての矢を保つものとする。
滑らかでない矢を個別に除く操作は用いない。

\begin{lemma}[generic flatnessと単位の滑らかさ]
\label{lem:generic-flat-full-groupoid}
$k$ を体，$V$ を整な有限型 $k$-スキームとし，
$\mathcal E\rightrightarrows V$ を有限型群oidスキームとする。
ある非空開集合 $V^0\subset V$ と整数 $g\ge0$ について，
$s:\mathcal E\to V$ は各 $v\in V^0$ の単位 $e(v)$ で滑らか，相対次元 $g$ と仮定する。
このとき非空稠密開集合 $W\subset V^0$ があり，全制限
\[
 \mathcal E_W:=\mathcal E\times_{V\times V}(W\times W)
       \underset{t}{\overset{s}{\rightrightarrows}}W
\]
の $s,t$ は全ての点で滑らか，相対次元 $g$ である。
$W$ は，あらかじめ指定した任意の非空開集合 $H\subset V$ の中に取れる。
\end{lemma}
\begin{proof}
整な基底 $V$ に対して $s:\mathcal E\to V$ にgeneric flatnessを適用する
\cite[Tag 052A]{StacksProject}。
ある非空開集合 $F\subset V$ 上で
$s:s^{-1}(F)\to F$ は平坦かつ有限表示となる。
$W\subset F\cap V^0\cap H$ を非空開集合に取る。
$V$ は整なのでこの共通部分は非空である。
対象の両端を $W$ に制限することは，基底変更と定義域の開制限の合成であるから，
$s:\mathcal E_W\to W$ は平坦かつ有限表示である。

幾何学的点 $x\in W(\Omega)$ と，そのsourceファイバーの幾何学的点
$\alpha:x\to y$ を取る。ここで $\Omega$ は代数的閉体とし，必要なら拡大する。
$y$ も $W(\Omega)$ に属する。
固定した $\alpha$ との右合成は $\Omega$-スキームの同型
\[
 R_\alpha:(\mathcal E_W)_{s=y}\xrightarrow{\ \sim\ }
                 (\mathcal E_W)_{s=x},
 \qquad h\longmapsto h\circ\alpha
\]
を与え，$e_y$ を $\alpha$ へ送る。
逆写像は $\alpha^{-1}$ との右合成である。
全制限を取っているため，この合成と逆写像はともに $\mathcal E_W$ の中に留まる。
単位での仮定により左辺のファイバーは $e_y$ で滑らか，局所次元 $g$ である。
従って右辺は $\alpha$ で滑らか，局所次元 $g$ である。
これを全幾何学的矢に適用すると，$s$ の全幾何学的ファイバーは滑らかで純次元 $g$ となる。
平坦かつ有限表示の射のファイバー判定
\cite[Tag 01V8]{StacksProject}から，$s$ は滑らかで相対次元 $g$ である。
逆射が $s,t$ を交換するため，$t$ にも同じ結論が従う。
\end{proof}

\begin{remark}\label{rem:flatness-before-fiber-translation}
平行移動で運んだのは幾何学的\emph{ファイバー}の滑らかさである。
射自体の滑らかさを得るための平坦性は，その前にgeneric flatnessで別に確保した。
この二段階を省いて，単位での滑らかさだけから全矢の滑らかさを結論してはいけない。
またgeneric flatnessの開集合 $F$ が可換零切断を含むとは限らない。
\end{remark}

\subsection{形式的普遍性だけでは零切断近傍を代数化できない}

本族の零切断近傍のDM表示には，規格化中心のCasimir Jacobianによる
全矢の持上げを用いた。以下の一般反例は，その追加の構造を使わず，
可換点の全Artin普遍性と有限表示対角だけから同じ結論を推論することを妨げる。

\begin{example}[消える安定化群とfullな商の障害]
\label{ex:formal-versality-not-open}
$V=\Spec\mathbb C[x,y]$，$z=(0,0)$，
$D=V(x)\cap D(y)$ と置く。
$V$ 上の群スキーム
\[
          G=V\amalg D\longrightarrow V
\]
を次のように定める：第一成分を単位とし，第二成分の元の平方を単位とする。
厳密には $D\times_VD=D$ を用いて第二成分同士の積を $D\hookrightarrow V$ とする。
他の積と逆射は自明に定まり，$D$ 上の群は $\mathbb Z/2$，$V\setminus D$ 上の群は自明である。
自明作用による群oid $G\rightrightarrows V$ を取る。

この群oidは次の性質を同時に持つ。
\begin{enumerate}
\item $(s,t):G\to V\times V$ はaffine，有限表示，quasi-finiteかつunramifiedである。
\item $s=t:G\to V$ は全ての単位で\'{e}taleである。
\item $z$ における全framed Artin変形群oidは
      $\operatorname{Hom}_{\mathrm{cont},\mathrm{loc}}(\mathbb C[[x,y]],-)$ という離散群oidである。
\item それにもかかわらず，$z$ を含むどの開集合 $W\subset V$ に対しても，
      全制限のfppf商 $[W/G_W]$ は代数スタックではない。
\end{enumerate}
\end{example}
\begin{proof}
$D\to V$ はaffineな有限表示の局所閉埋込みであり，
$G\to V$ の二成分は恒等射とこの局所閉埋込みである。
従って $G\to V$ はaffine，有限表示，quasi-finiteかつunramifiedであり，
対角 $V\hookrightarrow V\times V$ と合成しても同じ性質を持つ。
単位は開閉な第一成分にあるから(2)も従う。

局所Artin環 $A$ 上の射 $a:\Spec A\to V$ が $z$ に還元するなら，
$a^*y$ は冪零であるため $D\times_V\Spec A=\varnothing$ となる。
従って $G_A$ は自明な群スキームであり，そのtorsor群oidも自明である。
自明作用の商の対象は，$V$ への射とその引戻し群のtorsorであるから，
$z$ の形式近傍では $V$ への射そのものを離散に分類する。
これが(3)であり，接空間だけでなく全ての局所Artin環について成立する。

一方，$z\in\overline D$ なので，$z$ の任意の開近傍 $W$ は $D$ と交わる。
$D\cap W$ 上の第二成分では非零因子 $x$ が構造環を零化するため，
$G_W\to W$ は平坦ではない。
自明作用の商 $[W/G_W]$ が代数スタックであるための必要十分条件は，
$G_W\to W$ が平坦かつ局所有限表示であることに等しい
\cite[Tag 06PL]{StacksProject}。
従って(4)が従う。
\end{proof}

\begin{remark}[この反例が除外する推論]
非smoothな矢の像は $D$ であり，それ自身は $z$ に接しないが，その閉包は $z$ を含む。
従って矢のsmooth locusの逆像を単位で調べても，
$z$ を保った一つの基底開集合上で全矢をsmoothにすることはできない。
第二成分を捨てれば単位だけの群oidとなり商は $V$ になるが，
$D$ 上の自己同型を落とすため，元の全同値モジュライへの射は完全忠実ではない。
局所的な形式sliceや，限定した矢だけの商を構成しても，この問題は解消しない。

滑らかな基底との積を取れば，同じ反例を零切断全体に沿って作れる。
また $\mathbb A^g$ のpair群oid
$\mathbb A^g\times\mathbb A^g\rightrightarrows\mathbb A^g$ との積を取れば，
単位におけるsourceの相対次元を任意の $g\ge0$ にできる。
積の対象空間の開近傍の射影は $z$ を含む開集合であり，
pair群oidの全制限はその射影上の元の商と同値なので，非代数性の障害も残る。
これは本研究の特定の群oidが零切断近傍で非代数的であると主張する例ではない。
零切断での全Artin普遍性と既知の対角の性質だけから，
その近傍のfullな商の代数性を一般論として推論できないことを示す。
\end{remark}
